% -----------------------------------------------------------------------------
% Este código foi desenvolvido pelo Prof. Wesley Pacheco Calixto, em 26/11/2024.
% Este template foi criado para atender às necessidades acadêmicas dos discentes
% do  Instituto Federal de Educação, Ciência e Tecnologia de Goiás (IFG), unindo 
% precisão, padronização e automação no uso do LaTeX.
% -----------------------------------------------------------------------------
\begin{titlepage}
	\begin{center}
		\vspace*{2cm} % Espaço superior
		
			{DIOGO DOS REIS ALMEIDA \par} 
	        {HUAKSON HUILNNER SANTOS LIMA \par}
			\vspace{3cm}
		
		{\normalsize PROPOSTA DE ALGORITMO DE ROTEAMENTO EM REDES MESH DE BAIXO CUSTO COM ESP-NOW} \\[2cm] % Título do trabalho centralizado
	\end{center}
	
	% Texto formatado como citação direta
	\noindent
	\begin{flushleft}
		\setlength{\leftskip}{4cm} % Recuo à esquerda
		\setlength{\rightskip}{1cm} % Recuo à direita
		\setlength{\parindent}{0pt} % Remove a indentação
		Trabalho de Conclusão de Curso, no formato de monografia, apresentado ao curso de Bacharelado em Engenharia de Software, do Instituto Federal de Educação, Ciência e Tecnologia de Goiás - Campus Inhumas, como requisito parcial à obtenção do título de Bacharel em Engenharia de Software.
		
		\vspace{1cm}
		
		% Orientadores alinhados à esquerda
		\noindent
		Orientador: {KENYO ABADIO CROSARA FARIA} \\[0.5cm]
	\end{flushleft}
	
	\vfill
	\begin{center}
		% Local e data centralizados
		Inhumas \\[0.5cm]
		\today % Data automática
	\end{center}
	
\end{titlepage}
